% ECONOMICS_PROBLEM: {"id": "E52-44", "title": "Macroeconomic policy and information with dispersed beliefs", "jel_code": "E52", "source_id": "OP-0585"}
% !TeX program = xelatex
\documentclass[11pt,a4paper]{article}
\usepackage[margin=23mm]{geometry}
\usepackage{fontspec}
\setmainfont{Latin Modern Roman}
\usepackage{amsmath,amssymb,mathtools}
\usepackage{xcolor,enumitem,booktabs,longtable,array,xurl,hyperref}
\definecolor{muted}{HTML}{53636C}
\hypersetup{colorlinks=true,urlcolor=blue,linkcolor=blue}
\setlength{\parindent}{0pt}
\setlength{\parskip}{5pt}
\newcommand{\R}{\mathbb R}
\newcommand{\E}{\mathbb E}
\newcommand{\N}{\mathbb N}
\newcommand{\ind}{\mathbf 1}
\newcommand{\Var}{\operatorname{Var}}
\newcommand{\Cov}{\operatorname{Cov}}
\newcommand{\supp}{\operatorname{supp}}
\DeclareMathOperator*{\argmin}{arg\,min}
\DeclareMathOperator*{\argmax}{arg\,max}
\newcommand{\entrymeta}[3]{{\sffamily\small\color{muted}\textbf{\texttt{#1}}\quad|\quad #2\par #3\par}}
\newcommand{\Problemref}[1]{\texttt{#1}}
\newcommand{\JELref}[2]{\texttt{#2} (JEL \texttt{#1})}
\begin{document}
\section*{Macroeconomic policy and information with dispersed beliefs}
\textbf{Problem ID:} \texttt{E52-44}\quad \textbf{JEL:} \texttt{E52}\par
% BEGIN ECONOMICS STATEMENT
\label{problem:OP-0585}
\label{beyond:MA-04}
\entrymeta{MA-04}{Monetary policy and information}{Historical research agenda formalization}
\subsection*{Definitions and assumptions}
Use the infinite-horizon counterpart of the competitive economy with dispersed information in MA-03, a central-bank information filtration $\mathcal F_t^{\rm cb}$, admissible macroeconomic controls $a_t\in A$, and a family $\mathcal K$ of feasible public or private disclosure kernels. Controls must be $\mathcal F_t^{\rm cb}$-adapted; announcements and private messages are generated by the selected kernels and enter households' information sets. The instruments, government budget constraints and market-clearing effects are part of the primitives.

For each policy $(a,\kappa)$ let $\mathcal E(a,\kappa)$ be its nonempty competitive-equilibrium correspondence. Specify either a selection $e(a,\kappa)\in\mathcal E(a,\kappa)$ or an explicit welfare criterion over that correspondence. With a selected equilibrium, bounded measurable welfare $W$ and $\beta\in(0,1)$, define
\[
 J(a,\kappa)=\E^{e(a,\kappa)}\!\left[
 \sum_{t=0}^{\infty}\beta^tW(x_t,a_t,\mu_t)\right],\qquad
 V(\mathcal K)=\sup_{a,\ \kappa\in\mathcal K}J(a,\kappa),
\]
where $\mu_t$ is the equilibrium distribution of individual states and beliefs, rather than an exogenous input.
\subsection*{Formal question}
Under explicit policy, information and equilibrium-selection restrictions, prove attainment or characterize approximating optimal policies, derive dynamic programming or first-order conditions on a sufficient state, and determine how information affects welfare. Compare $V(\mathcal K_{\rm pub})$ and $V(\mathcal K_{\rm priv})$ for declared, equally costly disclosure technologies; characterize when transparency changes optimal stabilization or the allocation of information across agents.
\subsection*{Scope and status}
An arbitrary equilibrium correspondence makes the planner objective ambiguous unless selection is specified. More information need not have a universally signed welfare effect, and public versus private disclosure must be compared under a defined resource constraint. The problem is a sharpening of a historical policy agenda; it does not assert that these conclusions are unresolved for every model in the literature. Individual behavior remains competitive and the policy problem is a planner optimization.
\subsection*{References for the source of the problem}
N.~Gregory Mankiw and Ricardo Reis, \emph{Imperfect Information and Aggregate Supply}, NBER Working Paper~15773 (2010), \url{https://www.nber.org/system/files/working_papers/w15773/w15773.pdf}. Verified handbook version, Chapter~5 of \emph{Handbook of Monetary Economics} 3A: Section~5.2, ``Partial information and optimal transparency,'' and Section~7.3, ``Optimal policy with imperfect information.'' \url{https://personal.lse.ac.uk/reisr/papers/10-HofME.pdf}.

\subsection*{Common conventions: Source mathematical conventions}

\textbf{Well-defined objects.}
All probability spaces and measurable variables are specified or inherited
from a local statistical model. Expectations used as finite targets require
the stated integrability. A decision rule or estimator is measurable with
respect to its available information; randomized rules may use an auxiliary
independent random variable. Norms without a subscript are Euclidean in
finite-dimensional spaces. An optimizer is requested only when attainment
is assured; otherwise the question concerns the optimal value and
$\varepsilon$-optimal admissible rules for every $\varepsilon>0$.

\textbf{Identification.}
A structural model $m\in\mathcal M$ implies an observable law
$P_m^{\rm obs}$ and target $\theta(m)$. The target is point identified on
$\mathcal M$ if
\[
 P_{m_1}^{\rm obs}=P_{m_2}^{\rm obs}
 \quad\Longrightarrow\quad\theta(m_1)=\theta(m_2)
 \qquad(m_1,m_2\in\mathcal M).
\]
When this fails, the sharp identified set is
\[
 \Theta(P;\mathcal M)=\{\theta(m):m\in\mathcal M,
                                  P_m^{\rm obs}=P\}.
\]
Specifying an intervention or estimand does not establish identification.
Positivity, exclusion, causal sufficiency, measurement restrictions, and
transport assumptions are substantive local inputs, never automatic
consequences of notation.

\textbf{Minimax risk and nontrivial inference.}
For a statistical experiment with data law $P^{(n)}$, target $\theta(P)$,
and nonnegative loss $\ell$, define
\[
 \mathcal R_n(\mathcal P,\ell)
   =\inf_{\widehat\theta_n}\sup_{P\in\mathcal P}
      \E_{P^{(n)}}\ell(\widehat\theta_n,\theta(P)).
\]
The infimum is over measurable estimators. A sharp rate is a sequence
$r_n$ for which $0<c\leq C<\infty$, independent of $n$, satisfies
$c r_n\leq\mathcal R_n\leq C r_n$ for all sufficiently large $n$.
Squared-error parametric risk is of order $n^{-1}$; the corresponding
root-mean-squared error is of order $n^{-1/2}$. These different conventions
are distinguished locally. A uniform asymptotic confidence region $C_n$
must satisfy
\[
 \liminf_{n\to\infty}\inf_{P\in\mathcal P}
     \Pr_{P^{(n)}}\{\theta(P)\in C_n\}\geq 1-\alpha,
 \qquad 0<\alpha<1,
\]
together with an informative diameter or expected-width requirement.
Coverage alone permits the vacuous whole parameter space. Testing questions
similarly impose both size control and power against a declared separated
alternative class.

\textbf{Binary causal baseline.}
When an entry explicitly invokes this baseline, one observes independent
copies of $O=(X,W,Y)$, with $W\in\{0,1\}$ and potential outcomes $Y(0),Y(1)$.
Assume consistency $Y=Y(W)$, no interference, conditional exchangeability
$(Y(0),Y(1))\perp W\mid X$, and overlap
$\epsilon\leq e(X)\leq1-\epsilon$ almost surely for a fixed
$0<\epsilon<1/2$, where
\[
 e(x)=\Pr(W=1\mid X=x),\qquad
 m_w(x)=\E[Y\mid W=w,X=x],\qquad
 \tau(x)=m_1(x)-m_0(x).
\]
These assumptions identify conditional mean effects almost everywhere
under the covariate law. Pointwise or uniform effects require appropriate
regularity and a specified support. Continuous treatment, longitudinal
data, adaptive experiments, and network interference require their own
assumptions in place of this baseline. Bounded outcomes, densities, and
smoothness are additional assumptions only where stated.

\textbf{Function classes.}
On $[0,1]^d$, $\mathcal H_d^s(L)$ is the isotropic Hölder ball of radius
$L>0$: put $k=\lceil s\rceil-1$ and $\gamma=s-k\in(0,1]$, bound derivatives
through order $k$, and require the order-$k$ derivatives to be
$\gamma$-Hölder with constant at most $L$. Thus integer orders use a
Lipschitz final derivative convention. The class
$\operatorname{BV}_d(L)$ consists of functions $f\in L^1((0,1)^d)$ whose
distributional derivative is a finite vector measure and for which
$\|f\|_{L^1}+|Df|((0,1)^d)\leq L$.
An $L^\infty$ bound or a particular pointwise version is imposed separately
when needed. Sparse and neural-network classes require a declared
dictionary or architecture and coefficient bounds.

An anisotropic H\"older class with declared block smoothness indices means
coordinatewise H\"older regularity in each block, uniformly in the other
blocks, with all corresponding derivative and H\"older bounds at most
the stated radius. Derivative orders and the integer-order convention in
each block are those just given. Mixed-derivative bounds are additional
restrictions only when explicitly stated.

\textbf{Decision and computational questions.}
Policy regret compares admissible feasible policies under the same law and
constraints. In computational entries, finite data and thresholds have
rational encodings; running time is measured in their total bit length.
Real-valued equilibrium search uses the explicitly declared algebraic or
FIXP encoding. An approximation ratio for maximization is written
$\operatorname{OPT}/\operatorname{ALG}$ and for minimization as
$\operatorname{ALG}/\operatorname{OPT}$, with zero optima and infeasible
instances handled locally. Historical approximation bounds are not
automatically current bounds.
% END ECONOMICS STATEMENT
\end{document}
